\section{Conclusion}

We set out to test how far MPI's design transfers to multi-agent systems
built on language models, and the answer is: nearly all of it, and the
places where it does not are the interesting ones.

What transfers is the load-bearing structure. Communicators with isolated
contexts solve the library-composition problem for agent harnesses exactly
as they solved it for parallel libraries in 1993, and multi-agent frameworks
today exhibit the same silent interference that motivated them. The
collective catalogue, the matching rules, the group algebra, the epoch
structure of one-sided operations, two-phase collective I/O, the
revoke/shrink/agree recovery vocabulary, and the name-shifted profiling
interface all carry over with their semantics intact.

What does not transfer follows from four properties of the executor.
Context is consumed rather than occupied, which turns a receive into an
irreversible allocation, makes collective \emph{feasibility} a question that
must be answered before cost, and---as we found the hard way---means a
reduction tree must be planned against the root's cumulative ingest rather
than its per-round ingest. Reduction operators are semantic, so they must
declare their algebra and the runtime must refuse the schedules that algebra
does not license; the sharpest instance is that a commutative merge cannot
express precedence and is therefore unsound for any collective whose purpose
is agreement. Failure is a lattice rather than a binary, requiring
independent liveness and progress signals and a content-level error class.
And because computation costs five to seven orders of magnitude more than
communication, latency-optimal schedules win everywhere, which inverts
several of MPI's standard algorithm choices.

Our evaluation is small in scale and honest about it. Twelve agents
translating a book, eight building a query engine, and a fault-injection
study are not a proof that this design scales to a hundred; the
microbenchmarks that do reach 128 ranks measure the protocol rather than the
agents. But the runs were real, they surfaced failures we did not
anticipate---a launcher concurrency limit that deadlocked a collective, two
agents that skipped protocol steps and reasoned confidently from inferred
inputs---and each of those failures had a name and a remedy in the
distributed-systems literature, which is the point.

The broader claim we would defend is a negative one. The current generation
of agent protocols standardises the relationship between two parties, and
the current generation of agent frameworks each rebuilds group semantics
privately and incompatibly. That is the situation distributed-memory
computing was in before 1994, and it was resolved not by a better framework
but by a specification thin enough to implement many ways and precise enough
to build libraries on. We think the same move is available here, and we have
tried to show what it would look like.

\section*{Availability}

The protocol specification, reference implementation, test suite,
experiment harnesses, raw run data and the source of this paper are in the
accompanying repository.
